% !TeX root = ../sections/02_exercises.tex

\begin{exercise}
	R-9. 環 $A=\mathbb{Z}[\sqrt{2},\sqrt{3}]$ のイデアル
	\[
	I=(\sqrt{6})
	\]
	は $\sqrt{6}$ で生成された単項イデアルとする。
	\begin{enumerate}
		\item $\mathbb{Z}$-加群として $I$ は自由加群であることを示せ。
		また，基底を一組求めよ。
		\item $\mathbb{Z}$-加群として $A/I$ はねじれ加群で，
		\[
		A/I\cong
		\mathbb{Z}/6\mathbb{Z}\oplus
		\mathbb{Z}/6\mathbb{Z}
		\]
		に同型，すなわち不変系は $(6,6)$ であることを示せ。
	\end{enumerate}
\end{exercise}

\begin{background}
	$A=\mathbb{Z}[\sqrt{2},\sqrt{3}]$ は，
	$\mathbb{Z}$-加群として
	\[
	A=
	\mathbb{Z}
	\oplus
	\mathbb{Z}\sqrt{2}
	\oplus
	\mathbb{Z}\sqrt{3}
	\oplus
	\mathbb{Z}\sqrt{6}
	\]
	である。
	
	すなわち，任意の元は一意的に
	\[
	a+b\sqrt{2}+c\sqrt{3}+d\sqrt{6}
	\qquad
	(a,b,c,d\in\mathbb{Z})
	\]
	と表される。
	
	イデアル
	\[
	I=(\sqrt{6})
	\]
	は，
	\[
	I=\{\sqrt{6}x\mid x\in A\}
	\]
	で定義される。
	
	したがって，$A$ の一般元に $\sqrt{6}$ を掛けることで，
	$I$ の $\mathbb{Z}$-加群としての生成元が得られる。
\end{background}

\begin{strategy}
	任意の元 $x\in A$ を
	\[
	x=a+b\sqrt{2}+c\sqrt{3}+d\sqrt{6}
	\qquad
	(a,b,c,d\in\mathbb{Z})
	\]
	と書く。
	このとき
	\[
	\sqrt{6}x
	=
	a\sqrt{6}
	+b\sqrt{12}
	+c\sqrt{18}
	+d\sqrt{36}.
	\]
	ここで
	\[
	\sqrt{12}=2\sqrt{3},
	\qquad
	\sqrt{18}=3\sqrt{2},
	\qquad
	\sqrt{36}=6
	\]
	であるから，
	\[
	\sqrt{6}x
	=
	6d+3c\sqrt{2}+2b\sqrt{3}+a\sqrt{6}.
	\]
	
	したがって
	\[
	I=
	6\mathbb{Z}
	\oplus
	3\mathbb{Z}\sqrt{2}
	\oplus
	2\mathbb{Z}\sqrt{3}
	\oplus
	\mathbb{Z}\sqrt{6}.
	\]
	よって，$I$ は $\mathbb{Z}$-自由加群であり，一組の基底は
	\[
	6,\quad 3\sqrt{2},\quad 2\sqrt{3},\quad \sqrt{6}
	\]
	である。
	
	商加群 $A/I$ は，基底
	\[
	1,\sqrt{2},\sqrt{3},\sqrt{6}
	\]
	に関して
	\[
	\mathbb{Z}/6\mathbb{Z}
	\oplus
	\mathbb{Z}/3\mathbb{Z}
	\oplus
	\mathbb{Z}/2\mathbb{Z}
	\oplus
	0
	\]
	と見られる。
	
	さらに
	\[
	\mathbb{Z}/3\mathbb{Z}
	\oplus
	\mathbb{Z}/2\mathbb{Z}
	\cong
	\mathbb{Z}/6\mathbb{Z}
	\]
	であるから，
	\[
	A/I
	\cong
	\mathbb{Z}/6\mathbb{Z}
	\oplus
	\mathbb{Z}/6\mathbb{Z}
	\]
	となる。
\end{strategy}